% ===================================================================
%  அட்டவணை எண் 7 — குளிர்காலத்தில் கிராமம். வசந்த காலத்தில் பண்ணை வீடு.
%
%  Traduction, faite en 2026, en tamoul d'aujourd'hui, sur l'ido de
%  texto/io. Regles de la colonne : voir 10-attavanai-01.tex. Deux
%  scenes, deux numerotations : les cles portent c1 et c2.
%
%  LE MARRON CHAUD (42), ET LE TAMOUL TOMBE DU COTE DES ABSENTS. Ce
%  renvoi est le trou le mieux documente du releve : le telougou,
%  le marathe et l'egyptien n'ont pas le mot parce que le
%  chataignier ne pousse pas chez eux. Le coreen l'avait, et avec le
%  mot toute la scene — 군밤, les marrons grilles vendus dans la rue
%  en hiver. Le tamoul ne l'a pas davantage que le telougou : le
%  chataignier ne pousse pas au sud de l'Inde, et la colonne ecrit
%  செஸ்நட், l'emprunt. Quatre colonnes sans, une avec ; la coupure
%  suit la carte et non la famille de langues, et c'est la cinquieme
%  fois que ce releve le mesure.
%
%  LA BLAGUE DU VIEUX CORDONNIER, ET LE TROISIEME REFUS DU MEME
%  GENRE. « En ville j'etais shuifisto et je n'avais pas de clients ;
%  ici je suis shureparisto et le travail ne manque pas. » Le tamoul
%  a un nom de caste pour ce metier, et c'est aujourd'hui une injure
%  grave — la colonne ne l'ecrit pas plus que la colonne coreenne
%  n'avait ecrit 갖바치 ni la colonne telougoue కసాయి. Elle ecrit la
%  fonction, deux fois : செருப்பு செய்பவன் et செருப்பு பழுது
%  பார்ப்பவன். La blague tient toujours, parce qu'elle repose sur le
%  registre et non sur le mot.
%
%  LA FERME TAMOULE EST ENTIERE EN TAMOUL, et c'est la deuxieme fois
%  apres la cuisine du tableau 6 : கலப்பை la charrue, உழவுச் சால் le
%  sillon, பரம்பு la herse, மம்பட்டி la pioche, ராட்டை le rouet,
%  கிணறு le puits, தொழுவம் l'etable, மந்தை le troupeau, எரு le
%  fumier, கொண்டை la crete, மயில் le paon. Les betes ont toutes leur
%  nom et le tamoul les separe par un mot entier la ou l'ido met -ulo
%  et -ino : காளை le taureau contre பசு la vache, சேவல் le coq contre
%  கோழி la poule.
%
%  UN SEUL OISEAU MANQUE, ET C'EST L'OIE (69). Le tamoul dit வாத்து
%  pour le canard et n'a pas de mot courant pour l'oie, qui n'est pas
%  un oiseau de basse-cour au pays tamoul. La colonne ecrit
%  பெருவாத்து, le grand canard, et le declare ici plutot que de le
%  masquer sous வாத்து — ce qui aurait confondu les renvois (66),
%  (67) et (69) que le fac-simile distingue.
%
%  ET UNE PHRASE DE BROUILLON EST RESTEE DANS LE TEXTE, POUR LA
%  DEUXIEME FOIS. Au bloc c2-02-1 j'avais ecrit la phrase, vu qu'elle
%  sortait a l'envers, et ecrit la correction A LA SUITE au lieu de
%  remplacer : « ... — non, ce qu'il faut dire est ceci : ... ». Le
%  meme accident qu'au bloc 01-49 du tableau 5. renvoji.py le releve
%  a coup sur, parce que la phrase doublee double les renvois — mais
%  il ne le releverait PAS si le brouillon ne portait aucun renvoi.
%  La parade n'est donc pas dans l'outil : c'est de remplacer et non
%  d'ajouter, et de relire le bloc entier avant de passer au suivant.
%
%  TRENTE-DEUX PAIRES, DIX-NEUF RETOURNEMENTS. Le tableau se partage
%  en deux et le partage se lit dans le compte : la premiere scene
%  raconte des visites et enumere des artisans — quatre
%  retournements sur sept paires ; la seconde POSE la basse-cour
%  objet par objet — quinze sur vingt-cinq. C'est la seconde qui
%  coute, exactement comme la colonne coreenne l'avait dit.
% ===================================================================

%%K t07-noto-1 noto
\VUnotes{143.2mm}{%
(*) உணவு பற்றிய விவரங்களுக்கு அட்டவணை 6 ஐயும் 13 ஐயும் பார்க்க.%
}
%%K t07-apar-1 apar
\VUsaut{4.0mm}
\VUcentre{13.2pt}{90}{இரண்டாம் தொடர்}
\VUsaut{3.0mm}
\VUfilet{16mm}
\VUsaut{5.6mm}
\VUtitre{15.0pt}{60}{அட்டவணை எண் 7}
\VUsaut{3.0mm}

%%K t07-tit-1 sub
\VUcentre{12.0pt}{0}{\textit{முதல் காட்சி.}}
\VUsaut{3.0mm}

%%K t07-tit-2 sub
\VUpk{10.2pt}{40}{குளிர்காலத்தில் கிராமம்.}
\VUsaut{4.0mm}

%%K t07-c1-01-1 p
1. --- புத்தாண்டு விடுமுறையில் நான் என் தந்தையுடன் நோவூர்பொ என்னும் சிறு
கிராமத்திற்குப் போனேன் ; அங்கே அவருக்குச் சில வணிக வேலைகள் இருந்தன.

%%K t07-c1-02-1 p
2. --- நகரத்திற்கும் அந்தச் சிற்றூருக்கும் இடையே ஓடும் \VUgras{அஞ்சல்
வண்டியில்} \textsuperscript{(11)} போனோம் ; ஆனால் மிகவும் குளிராக
இருந்ததால், வசதி குறைந்த அந்த வண்டியில் பயணம் மிகவும் இனிமையாக
இல்லை.

%%K t07-c1-03-1 p
3. --- அதனால், \VUgras{வண்டிக்காரர்} தன் வண்டியை \textit{வெள்ளைக் கரடி}
என்னும் \VUgras{தங்கும் விடுதிக்கு} \textsuperscript{(2)} முன்னால்
சதுக்கத்தில் நிறுத்தியதைப் பார்த்த போது நமக்கு உண்மையான மகிழ்ச்சி
ஏற்பட்டது. \VUgras{விடுதிக்காரர்} \textsuperscript{(4)} தன் வீட்டு
வாசற்படியில் நின்று அமைதியாகத் தன் \VUgras{புகைக் குழலைப்}
புகைத்துக் கொண்டு நமக்குக் காத்திருந்தார்.

%%K t07-c1-03-2 p
அவர் நம்மை உள்ளே அழைத்தார் ; அடுப்பில் படபடத்த கொடிக் குச்சி
நெருப்புக்கு முன்னால் அமர்வதற்கு நான் இரண்டாம் முறை சொல்லக்
காத்திருக்கவில்லை. அப்போது பழைய \VUgras{பெயர்ப் பலகையைக்}
\textsuperscript{(3)} கிறீச்சிட வைத்த கடுமையான \VUgras{வட காற்றை} நான்
மிகவும் கேலி செய்தேன் ; அது புகைபோக்கியிலிருந்து வெளியேறும்
\VUgras{புகையைக்} \textsuperscript{(1)} கருஞ் சுழல்களாக அடித்துச்
சென்றது.

%%K t07-c1-04-1 p
4. --- காலை உணவின் நேரம் அது. மேலும் காத்திராமல், நமக்குப் பரிமாறிய
எளிமையான ஆனால் சுவையான உணவை நன்றாகச் சாப்பிட்டோம் (*).

%%K t07-tit-3 sub
\VUpk{10.2pt}{40}{சில வருகைகள்.}
\VUsaut{3.0mm}

%%K t07-c1-05-1 p
5. --- காலை உணவுக்குப் பிறகு, காப்பீட்டு முகவரான என் தந்தையுடன் அவருடைய
வாடிக்கையாளர்களிடம் போனேன். முதலில் \VUgras{கொல்லரைப்}
\textsuperscript{(5)} பார்த்தோம் ; அவர் விடுதிக்கு அருகில் குடியிருக்கிறார்.
அவர் ஒரு குதிரைக்கு லாடம் அடிப்பதில் ஈடுபட்டிருந்தார். அவருடைய உதவியாளின்
வலிமையை நான் வியந்தேன் ; அவர் தன் தோல் முன்னாடையின் மேல் குதிரையின்
\VUgras{குளம்பை} உறுதியாகப் பிடித்திருந்தார் ; அதே வேளையில் கொல்லர்
வலிமையான சுத்தி அடிகளால் \VUgras{லாடத்தைப்} \textsuperscript{(6)}
பொருத்தினார்.

%%K t07-c1-06-1 p
6. --- பிறகு கிராமத்தின் \VUgras{செருப்புத் தைப்பவரிடம்}
\textsuperscript{(7)}, பழைய யாகோபுஸிடம் போனோம். அவருடைய பெயர்ப் பலகையாக
மிக அழகான \VUgras{பூட்ஸ்} இருந்தாலும், அவர் ஓட்டை விழுந்த ஒரு ஜோடி
பழைய செருப்புகளுக்குப் புதிய அடி போடுவதிலேயே திருப்தி அடைந்தார்.

%%K t07-c1-06-2 p
அந்த முதியவர் சிரித்துக் கொண்டே நம்மிடம் சொன்னார் : « நகரத்தில் நான்
« செருப்பு செய்பவன் » ; எனக்கு வாடிக்கையாளர்கள் இல்லை. இங்கே நான்
« செருப்பு பழுது பார்ப்பவன் » ; வேலை நிறையக் கிடைக்கிறது. »

%%K t07-c1-07-1 p
7. --- அவர் தன் அயலவரான \VUgras{முடி திருத்துபவரைப்} \textsuperscript{(8)}
பற்றிச் சொன்னார் ; அவருடைய வாடிக்கையாளர்கள் மனநிறைவாக இல்லை. அவருடைய
\VUgras{சவரக் கத்திகள்} எப்போதும் மொட்டையாக இருக்கின்றன ; அவருடைய
\VUgras{கத்தரிக்கோல்கள்} ஒரு போதும் நன்றாக வெட்டுவதில்லை.

ஆனால் அவர்தான் கிராமத்தின் ஒரே முடி திருத்துபவர் என்பதால், அவரிடமே
சவரம் செய்து கொள்ளவும் தலை முடி வெட்டிக் கொள்ளவும் கட்டாயம் ஆகிறது.
மேலும் அவர் கஞ்சரும் இனிமையற்றவரும் ஆவார்.

%%K t07-c1-08-1 p
8. --- நல்ல வேளையாக மற்ற \VUgras{அயலவர்கள்} அவரைப் போல் இல்லை.
எடுத்துக்காட்டாக \VUgras{வண்டி செய்பவர்} \textsuperscript{(9)} நல்ல
மனிதர், உழைப்பாளி, உதவி செய்பவர். அவரிடம் வண்டியைப் பழுது பார்ப்பது
ஒரு மகிழ்ச்சி. அவருடைய வேலை எப்போதும் நேர்த்தியாக முடிந்திருக்கும்.

%%K t07-c1-09-1 p
9. --- \VUgras{சேணம் செய்பவரைப்} \textsuperscript{(10)} பற்றியும் பழைய
யாகோபுஸ் முறையிடவில்லை ; வேலை அவசரமாக இருக்கும் போது
\VUgras{சேணங்களைத்} தைக்க அவருக்குச் சில வேளைகளில் உதவுகிறார்.

%%K t07-c1-09-2 p
என் தந்தை அவரிடம் அவருடைய குடும்பத்தைப் பற்றிக் கேட்டார் : « எல்லாரும்
நலமாக இருக்கிறார்கள். என் பேத்தி \VUgras{பள்ளியில்} \textsuperscript{(12)}
இருக்கிறாள். அவளுடைய \VUgras{ஆசிரியை} \textsuperscript{(13)} அவளைப் பற்றி
மிகவும் மனநிறைவாக இருக்கிறாள் ; அவள் நல்ல \VUgras{மாணவி}
\textsuperscript{(14)}. என் கெட்ட மகனைப் போல் அவள் இல்லை ; அவன்
விளையாடுவதைப் பற்றியே நினைக்கிறான். \VUgras{ஆசிரியரால்}
\textsuperscript{(15)} அவனுக்கு எதையும் கற்பிக்க முடியவில்லை. »

%%K t07-tit-4 sub
\VUsaut{3.0mm}
\VUpk{10.2pt}{40}{கிராமத்து அளவளாவல்கள்.}
\VUsaut{3.0mm}

%%K t07-c1-10-1 p
10. --- பிறகு அந்த முதியவர் கிராமத்துச் செய்திகளை நமக்குச் சொன்னார்.
புதிய பள்ளி ஒன்று கட்டப் போகிறார்கள் ; ஏனெனில் \VUgras{ஊராட்சி
மன்றத்தில்} அமைந்திருந்த பள்ளி மிகவும் சிறியதாகி விட்டது. அது எளிதுதான் ;
ஏனெனில் \VUgras{மாவு ஆலைக்காரரின்} தோட்டத்தில் ஒரு பகுதியை
வாங்கினால் போதும். அந்தப் பரிதாபமான மனிதருக்கு அது மிகவும் பயன் தரும்.
குளிரால் \VUgras{மாவு ஆலையின்} \textsuperscript{(16)} \VUgras{திரிகைக்
கற்களால்} இனிக் கோதுமையை அரைக்க முடிவதில்லை ; ஏனெனில்
\VUgras{சக்கரத்தை} \textsuperscript{(17)} \VUgras{பனிக்கட்டி}
\textsuperscript{(25)} உறைய வைத்து நிறுத்தி விட்டது ; அது
\VUgras{ஓடையின்} \textsuperscript{(23)} பனிக்கட்டி.

%%K t07-c1-11-1 p
11. --- « கட்டிடங்களைப் பற்றிச் சொன்னால், நம் புதிய \VUgras{தேவாலயத்தை}
\textsuperscript{(18)} நீங்கள் பார்த்தீர்களா? பனிப் போர்வையின் கீழ் அது
மிகவும் அழகாக இருக்கிறது. \VUgras{காகங்கள்} \textsuperscript{(19)} தங்கள்
பழைய மணிக் கோபுரத்தை இனி அடையாளம் காணாமல், \VUgras{கல்லறைத்
தோட்டத்தின்} \textsuperscript{(20)} பெரிய மரத்தின் மேல் சோகமாக
அமர்ந்திருக்கின்றன. »

%%K t07-c1-12-1 p
12. --- கல்லறைத் தோட்டத்தைப் பற்றிய அந்த எண்ணம், என் தந்தைக்குத்
தெரிந்தவர்களுள் கடந்த ஆண்டு இறந்தவர்களைச் செருப்புத் தைப்பவருக்கு
நினைவூட்டியது. « எத்தனை \VUgras{கல்லறைகள்} \textsuperscript{(21)}, என்
நல்ல ஐயா, \VUgras{சைப்ரஸ் மரத்தின்} \textsuperscript{(22)} கீழ்
மற்றவற்றுடன் சேர்ந்தன! நம் பழைய நோட்டரியை மரணம் அறுத்து விட்டது.
கிராமத்தார் அனைவரும் அவருடைய \VUgras{இறுதிச் சடங்கில்} கலந்து
கொண்டார்கள். »

%%K t07-c1-13-1 p
13. --- « இதோ அவருடைய இடத்திற்கு வந்தவர் ; ஓடையின் மேல் சறுக்கி
விளையாடுகிறார். உடைந்து போன தன் \VUgras{பனிச் சறுக்கு காலணியின்}
\textsuperscript{(26)} வாரை என்னிடம் பழுது பார்க்கச் சொல்ல இப்போதுதான்
வந்தார். »

%%K t07-c1-14-1 p
14. --- « ஓடையின் கரையில் இதோ புதிய \VUgras{ஊர்க் காவலரும்}
\textsuperscript{(27)}. அவர் முன்னாள் காவலர் ; கிராமத்துக் குறும்புச்
சிறுவர்களை அச்சுறுத்துகிறார். அவருடைய \VUgras{கால் உறைகளையும்}
\textsuperscript{(29)} அவருடைய \VUgras{காவல் தொப்பியையும்}
\textsuperscript{(28)} பார்த்தவுடன் அவர்கள் ஓடி விடுகிறார்கள். »

%%K t07-c1-15-1 p
15. --- « ஆ! இதோ பழைய இயோசெஃப் ; ரொட்டி சுடுபவருக்காக \VUgras{விறகுக்
கட்டுகள் ஏற்றிய சிறு வண்டியை} \textsuperscript{(30)} ஓட்டி வருகிறார்.
--- உண்மைதான், என்றார் என் தந்தை ; அவருடைய \VUgras{கோவேறு கழுதைகளின்}
\textsuperscript{(31)} பூட்டை நான் அடையாளம் காண்கிறேன். அவரிடம் பேச
வேண்டியிருக்கிறது ; இந்த வாய்ப்பைப் பயன்படுத்திக் கொள்கிறேன்.
மீண்டும் சந்திப்போம், யாகோபுஸ் அப்பா. »

%%K t07-c1-16-1 p
16. --- நாம் வெளியே வந்த அந்த நேரத்தில், கிராமத்துச் சிறுவர்கள்
பள்ளியை விட்டு வெளியேறி, ஓடிக் கொண்டு \VUgras{சறுக்கும் இடத்திற்கு}
\textsuperscript{(33)} விரைந்தார்கள் — \VUgras{விழுந்து}
\textsuperscript{(34)} கை கால் உடையும் அபாயத்தையும் பொருட்படுத்தாமல்.
அவர்களுள் ஒருவன் வண்டி செய்பவரிடமிருந்து தன் \VUgras{சறுக்கு வண்டியை}
\textsuperscript{(32)} எடுத்து, அதில் தன் தோழர்களுள் ஒருவனை உட்கார
வைத்து, ஓடிக் கொண்டு புறப்பட்டான்.

%%K t07-c1-17-1 p
17. --- வேறு சிலர் \VUgras{பனிக் குவியலுக்கு} \textsuperscript{(37)}
விரைந்தார்கள் ; அது \VUgras{பாலத்திற்கு} \textsuperscript{(40)} அருகில்
இருக்கிறது. நேற்று அவர்கள் எழுப்பிய \VUgras{பனிச் சிலையை}
\textsuperscript{(39)} அவர்கள் தாக்கத் தொடங்கினார்கள் ; அதற்கு அவர்கள்
தொப்பியும், புகைக் குழலும், தோளில் தொங்கும் பள்ளிப் பையும்
போட்டிருந்தார்கள்.

%%K t07-c1-18-1 p
18. --- உறைய வைக்கும் காற்றையும் குளிரையும் அவர்கள் பெரிதாகப்
பொருட்படுத்துவதில்லை. அவர்களுள் பெரும்பாலோரிடம் \VUgras{கழுத்துக்
கம்பளிகூட} \textsuperscript{(35)} இல்லை. \VUgras{பனி உருண்டைச்}
\textsuperscript{(38)} சண்டைக்குப் பிறகு மீண்டும் கதகதப்பாக, மிகவும்
சூடான ஒரு கை நிறைய \VUgras{செஸ்நட் பழங்கள்} \textsuperscript{(42)}
அவர்களுக்குப் போதும் ; அவற்றை அவர்கள் பழைய \VUgras{வியாபாரி}
யாகோபினாவிடம் வாங்குகிறார்கள் ; அவள் மிகவும் மெல்லிய தன்
\VUgras{கிழிந்த சால்வையின்} \textsuperscript{(43)} கீழ் குளிரால்
நடுங்குகிறாள்.

%%K t07-tit-5 sub
\VUsaut{3.0mm}
\VUcentre{12.0pt}{0}{\textit{இரண்டாம் காட்சி.}}
\VUsaut{3.0mm}

%%K t07-tit-6 sub
\VUpk{10.2pt}{40}{வசந்த காலத்தில் பண்ணை வீடு.}
\VUsaut{4.0mm}

%%K t07-c2-01-1 p
1. --- குளிர்காலத்தின் எல்லா மகிழ்ச்சிகளையும் மீறி, \VUgras{வசந்த
காலத்தின்} வருகையை ஆழ்ந்த மகிழ்ச்சியுடன் வரவேற்கிறோம்.

%%K t07-c2-01-2 p
மே மாதத்தில், மாலை நேரத்தில், பண்ணை முற்றம் எத்தனை மகிழ்ச்சியான
காட்சியாகத் தெரிகிறது!

%%K t07-c2-02-1 p
2. --- அடிவானத்தில் \VUgras{சூரியன்} \textsuperscript{(1)} மெல்லப்
படுக்கிறது ; \VUgras{நாட்டுப்புறத்தை} \textsuperscript{(3)} தன் செந்நிற
\VUgras{ஒளிக் கதிர்களால்} \textsuperscript{(2)} நிரப்புகிறது. உழைப்பான
\VUgras{உழவன்} \textsuperscript{(6)} தன் \VUgras{வயலை}
\textsuperscript{(4)} உழ அவசரப்படுகிறான்.

%%K t07-c2-02-2 p
தன் \VUgras{கலப்பையால்} \textsuperscript{(5)} — அது வலிமையான
\VUgras{குதிரையுடன்} \textsuperscript{(7)} பூட்டப்பட்டிருக்கிறது —
அவன் \VUgras{உழவுச் சாலை} \textsuperscript{(8)} கீறுகிறான் ; அதில் \VUgras{விதைப்பவன்}
\textsuperscript{(9)} கை நிறைய \VUgras{விதையை} எறிவான் ; அது பொன்னிற
தானியக் கதிர்களைத் தரும்.

%%K t07-c2-03-1 p
3. --- \VUgras{புல் அறுப்பவர்கள்} \textsuperscript{(11)} தங்கள்
அரிவாள்களுடன் புல்வெளியின் புல்லை அறுத்தார்கள் ; உலர்த்துபவர்கள்
\VUgras{உலர் புல்லை} \textsuperscript{(10)} \VUgras{முள்ளுக் கவைகளாலும்}
\textsuperscript{(12)} இரும்புச் சீப்புகளாலும் புரட்டி உலர்த்தினார்கள் ;
இதோ அவர்கள் திரும்பி வருகிறார்கள்.

%%K t07-c2-03-2 p
சிலர் மாடுகள் இழுக்கும் பாரமான வண்டிக்கு முன்னால் நடக்கிறார்கள் ;
தங்கள் கருவியைத் தோளில் சுமந்து வருகிறார்கள். வேறு சிலர், இன்னும்
சோம்பலானவர்கள், மணக்கும் உலர் புல்லின் மேல் வசதியாக அமர்ந்து
கொண்டார்கள்.

%%K t07-c2-04-1 p
4. --- கடந்து போகும் போது அவர்கள் \VUgras{தோட்டக்காரருக்கு}
\textsuperscript{(16)} நட்பான வணக்கம் சொல்கிறார்கள் ; அவர்
\VUgras{பழத் தோட்டத்தில்} \textsuperscript{(13)} \VUgras{பழ மரங்களைக்}
கத்தரிப்பதிலும் \VUgras{பேரி மரங்களை} \textsuperscript{(14)}
ஒட்டுவதிலும் ஈடுபட்டிருக்கிறார். \VUgras{ப்ளம் மரங்கள்}
\textsuperscript{(15)} பூக்கத் தொடங்குகின்றன ; நன்றாக எரு போட்ட
\VUgras{காய்கறித் தோட்டத்தில்} \textsuperscript{(17)} காய்கறிகளும்
முட்டைக்கோசுகளும் கீரைகளும் ஏற்கெனவே நன்றாக இருக்கின்றன.

%%K t07-c2-04-2 p
சிறு சுவரின் மேல் அழகான \VUgras{மயில்} \textsuperscript{(18)} தன்
வாலை விரிக்கிறது ; தன் மிக அழகான இறகுகளை வியக்க வைக்கிறது.

%%K t07-tit-7 sub
\VUpk{10.2pt}{40}{பண்ணை முற்றம்.}
\VUsaut{3.0mm}

%%K t07-c2-05-1 p
5. --- \VUgras{குதிரை லாயத்தின்} \textsuperscript{(19)} கதவுகள்
திறந்திருக்கின்றன ; தீவனத் தொட்டியையும் \VUgras{படுக்கை வைக்கோலையும்}
\textsuperscript{(23)} பார்க்கிறோம் ; அவை \VUgras{பெண் குதிரையுடையவை}
\textsuperscript{(21)} ; அந்தக் குதிரையை \VUgras{குதிரைக் காரர்}
\textsuperscript{(20)} இப்போதுதான் அவிழ்த்து விட்டார். நீண்ட கால்களால்
மிகவும் நகைச்சுவையாகத் தெரியும் \VUgras{குட்டிக் குதிரை}
\textsuperscript{(22)} துள்ளிக் கொண்டு தன் தாயைப் பின்தொடர்கிறது.

%%K t07-c2-06-1 p
6. --- \VUgras{கிணற்றுக்கு} \textsuperscript{(29)} அருகில்
\VUgras{பசுக்கள்} \textsuperscript{(25)} \VUgras{மரத் தொட்டியில்}
\textsuperscript{(26)} நீர் குடிக்கப் போகின்றன ; அந்தத் தொட்டி
அவற்றுக்குத் தண்ணீர்த் தொட்டியாகப் பயன்படுகிறது.
\VUgras{கன்றுக்குட்டி} \textsuperscript{(24)} இந்த வாய்ப்பைப் பயன்படுத்திப்
பால் குடிக்கிறது ; \VUgras{காளை} \textsuperscript{(27)} தீவனக் களஞ்சியத்தின்
நல்ல மணத்தை முகர்ந்து மகிழ்ச்சியாகக் கத்துகிறது, வலிமையான
\VUgras{கொம்புகளுடன்} \textsuperscript{(28)} தன் தலையைப் பெருமையுடன்
உயர்த்துகிறது.

%%K t07-c2-07-1 p
7. --- எல்லாரும் வேலை செய்கிறார்கள். \VUgras{வேலைக்காரி}
\textsuperscript{(32)} கிணற்றின் \VUgras{கயிற்றைக்} \textsuperscript{(31)}
கையால் பிடித்திருக்கிறாள். கால்நடைகளுக்கு நீர் கொடுக்க
\VUgras{வாளியால்} \textsuperscript{(30)} நீர் இறைக்கிறாள். \VUgras{வைக்கோல்
களஞ்சியத்திற்கு} \textsuperscript{(33)} அருகில் ஒரு மனிதர்
வைக்கோலைச் சீவுகிறார் ; \VUgras{வீட்டின்} \textsuperscript{(34)}
கதவுக்கு முன்னால் பழைய \VUgras{நூல் நூற்பவள்} \textsuperscript{(35)}
தன் ராட்டையாலும் தன் \VUgras{கதிர்க் கோலாலும்} பழைய காலத்தைப் போலவே
\VUgras{ஆளியையோ} \VUgras{சணலையோ} நூற்கிறாள்.

%%K t07-tit-8 sub
\VUsaut{3.0mm}
\VUpk{10.2pt}{40}{பண்ணையாளர்.}
\VUsaut{3.0mm}

%%K t07-c2-08-1 p
8. --- \VUgras{பண்ணையாளர்} \textsuperscript{(36)} இந்த வேலைகள்
அனைத்தையும் நடத்துகிறார், தன் வேலையாட்களுக்குக் கட்டளைகள்
கொடுக்கிறார். \VUgras{பரம்பையும்} \textsuperscript{(40)}
\VUgras{மம்பட்டியையும்} \textsuperscript{(39)} கூரையின் கீழ் வைக்கச்
சொல்கிறார் ; அவற்றை \VUgras{கூட்டிற்கு} \textsuperscript{(38)} அருகில்
மறந்து விட்டார்கள் ; அது \VUgras{நாயினுடையது} \textsuperscript{(37)}.

%%K t07-c2-09-1 p
9. --- அவருடைய கூக்குரல்கள் ஒரு \VUgras{புறாவை} \textsuperscript{(41)}
பயமுறுத்துகின்றன ; அது சிறகடித்து \VUgras{புறாக் கூட்டிற்குள்}
\textsuperscript{(42)} பறக்கிறது. \VUgras{கோழிகளையும்}
\textsuperscript{(43)} பயமுறுத்துகின்றன ; அவை \VUgras{கோழிக் கூட்டிற்குள்}
\textsuperscript{(44)} அவசரமாகத் திரும்பி விடுகின்றன.

%%K t07-c2-10-1 p
10. --- \VUgras{ஆட்டுத் தொழுவத்திற்கு} \textsuperscript{(48)} முன்னால்,
இதோ பழைய \VUgras{ஆட்டு இடையன்} \textsuperscript{(45)} ; அவன் தன்
\VUgras{மந்தையைத்} \textsuperscript{(49)} திரும்ப அழைத்து வருகிறான்.
தன்னிடமிருந்து ஒரு போதும் பிரியாத தன் \VUgras{இடையன் கோலைப்}
\textsuperscript{(46)} பிடித்துக் கொண்டு — நிறம் மங்கிய தன்
\VUgras{மேற்சட்டையைப்} \textsuperscript{(47)} போலவே அதுவும் எப்போதும்
அவனுடன் இருக்கிறது —, \VUgras{செம்மறிக் கடாக்களையும்}
\textsuperscript{(50)}, \VUgras{செம்மறி ஆடுகளையும்} \textsuperscript{(53)},
\VUgras{ஆட்டுக் குட்டிகளையும்} \textsuperscript{(54)},
\VUgras{வெள்ளாடுகளையும்} \textsuperscript{(51)}, \VUgras{வெள்ளாட்டுக்
குட்டிகளையும்} \textsuperscript{(52)} தொழுவத்திற்கு அழைத்துச்
செல்கிறான். அவனுடைய நம்பிக்கையான \VUgras{நாய்} \textsuperscript{(55)}
ஆர்கஸ் கடைசியாக வருகிறது ; தன் குரைப்பால் பின்தங்குபவர்களை
விரைவுபடுத்துகிறது.

%%K t07-c2-11-1 p
11. --- இந்த ஆரவாரமும் அசைவும் நல்ல \VUgras{கறவைப் பசுவின்}
\textsuperscript{(56)} அமைதியைக் கலைக்கவில்லை ; அது தன்
\VUgras{மணியைக்} \textsuperscript{(57)} கிலுக்குகிறது ; அதே வேளையில்
\VUgras{தொழுவத்தின்} \textsuperscript{(58)} கதவுக்கு முன்னால் அதைக்
கறக்கிறார்கள்.

%%K t07-c2-12-1 p
12. தான் பால் கொடுக்க வேண்டும் என்பதை அது புரிந்து கொள்வது போல்
தெரிகிறது ; அப்போதுதான் \VUgras{பண்ணைத் தலைவி} \textsuperscript{(60)}
\VUgras{வெண்ணெய்த் தாழியில்} \textsuperscript{(61)} \VUgras{வெண்ணெயையோ},
மிகச் சிறந்த \VUgras{சீஸ்களையோ} \textsuperscript{(64)} தயார் செய்ய
முடியும் ; அந்தச் சீஸ்கள் \VUgras{பெரும் தொட்டியின்}
\textsuperscript{(63)} மேல் வைத்த பலகையிலிருந்து சொட்டுகின்றன.

%%K t07-c2-13-1 p
13. --- \VUgras{பால் அறை} \textsuperscript{(59)} முழுவதையும்
\VUgras{பண்ணை வேலையாள்} \textsuperscript{(65)} மிகவும் சுத்தமாக
வைத்திருக்கிறார் ; அவர் இப்போதுதான் \VUgras{பால் பானைகளைத்}
\textsuperscript{(62)} தான் கையில் சுமக்கும் பெரிய வாளியில்
கழுவினார்.

%%K t07-tit-9 sub
\VUsaut{3.0mm}
\VUpk{10.2pt}{40}{கோழிப் பண்ணை.}
\VUsaut{3.0mm}

%%K t07-c2-14-1 p
14. --- தன் தடித்த மரக் காலணிகளுடன் அவர் கொஞ்சம் கனமாக நடக்கிறார் ;
அதனால் \VUgras{வான்கோழியையும்} \textsuperscript{(68)}, \VUgras{பெண்
வாத்துகளையும்} \textsuperscript{(66)}, \VUgras{ஆண் வாத்துகளையும்}
\textsuperscript{(67)}, \VUgras{பெருவாத்துகளையும்} \textsuperscript{(69)}
விரட்டி விடுகிறார் ; அவை \VUgras{தங்கள் அலகை} \textsuperscript{(70)}
அகலமாகத் திறந்து கலக்கத்துடன் கத்துகின்றன.

%%K t07-c2-14-2 p
\VUgras{சேவல்} \textsuperscript{(71)}, இன்னும் சண்டைக்காரன், தன் சிவப்பு
\VUgras{கொண்டையை} \textsuperscript{(72)} நிமிர்த்துகிறது, தன்
\VUgras{வாலின்} \textsuperscript{(73)} இறகுகளை உதறுகிறது, முழங்கும்
« கொக்கரக்கோ » ஒலியை எழுப்புகிறது.

%%K t07-c2-15-1 p
15. --- \VUgras{தாய்க் கோழி} \textsuperscript{(74)} \VUgras{எருவின்}
\textsuperscript{(81)} மேல் தன் \VUgras{குஞ்சுகளுடன்}
\textsuperscript{(75)} கிளறித் தேடுகிறது. \VUgras{பன்றிக் குட்டி}
\textsuperscript{(78)} தன் தாயிடம் ஓடுகிறது — பன்றிக் கொட்டிலுக்கு
அருகில் நாம் பார்க்கும் அந்தப் பெரிய \VUgras{பெண் பன்றியிடம்}
\textsuperscript{(77)}.

%%K t07-c2-16-1 p
16. --- தங்கள் அறைகளில் \VUgras{முயல்கள்} \textsuperscript{(79)}
மென்மையான \VUgras{புல்லை} \textsuperscript{(80)} ஆவலுடன் தின்கின்றன ;
அதைப் பண்ணையாளரின் மகள் அவற்றுக்குக் கொண்டு வருகிறாள். இயோனீனா தன்
முயல்களை மிகவும் விரும்புகிறாள் ; நாள்தோறும் கோழிக் கூட்டில் புதிய
\VUgras{முட்டைகளை} \textsuperscript{(76)} எடுத்த பிறகு, தன் சிறு
நண்பர்களைப் பார்க்க வருகிறாள்.

\VUsaut{6.0mm}
\VUfilet{22mm}
